\documentclass[a4paper,12pt]{report}
\usepackage[margin=1in]{geometry}
\usepackage[T1]{fontenc}
\usepackage[utf8]{inputenc}
\usepackage{lmodern}
\usepackage{microtype}      % subtle font expansion and margin protrusion
\emergencystretch=3em   
\usepackage{tikz}
\usetikzlibrary{calc}
\usepackage{graphicx}
\usepackage{titlesec}
\usepackage{titletoc}              
\usepackage{fancyhdr}
\usepackage{amsmath}
\usepackage{listings}
\usepackage{xcolor}
\usepackage{caption}
\usepackage{subcaption}
\usepackage{array}
\usepackage{multirow}
\usepackage{booktabs}
\usepackage{makecell}
\usepackage{float}
\usepackage{lipsum}
\usepackage{adjustbox}
\usepackage{setspace}
\usepackage{indentfirst}
\setstretch{1.2}
\usepackage[colorlinks=true, linkcolor=black, citecolor=black, urlcolor=black]{hyperref}
\usepackage[backend=biber, style=numeric, sorting=none]{biblatex}
\addbibresource{references.bib}

% vietnamese
\usepackage[T5]{fontenc}
\usepackage[utf8]{inputenc}
\usepackage{lmodern}

\setlength{\headheight}{14.5pt}
\pagestyle{fancy}
\lhead{S\&T Management}
% Right header: "IV. Title" instead of "CHAPTER IV. TITLE", so it no longer overlaps the left header
\renewcommand{\chaptermark}[1]{\markboth{\thechapter.\ #1}{}}

\definecolor{codegreen}{rgb}{0,0.6,0}
\definecolor{codegray}{rgb}{0.5,0.5,0.5}
\definecolor{codepurple}{rgb}{0.58,0,0.82}
\definecolor{backcolour}{rgb}{0.95,0.95,0.92}
\lstdefinestyle{mystyle}{
    backgroundcolor=\color{backcolour},
    commentstyle=\color{codegreen},
    keywordstyle=\color{magenta},
    numberstyle=\tiny\color{codegray},
    stringstyle=\color{codepurple},
    basicstyle=\ttfamily\footnotesize,
    breaklines=true,
    captionpos=b,
    keepspaces=true,
    numbers=left,
    numbersep=5pt,
    showspaces=false,
    showstringspaces=false,
    showtabs=false,
    tabsize=2
}
\lstset{style=mystyle}

% ----- Đánh số kiểu I. / 1. / 2.1. -----
\renewcommand{\thechapter}{\Roman{chapter}}
\renewcommand{\thesection}{\arabic{section}}
\renewcommand{\thesubsection}{\thesection.\arabic{subsection}}

\titleformat{\chapter}[block]
  {\LARGE\bfseries}{\thechapter.}{0.5em}{\MakeUppercase}
\titlespacing*{\chapter}{0pt}{0pt}{1cm}
\titleformat{\section}[block]
  {\large\bfseries}{\thesection.}{0.5em}{}
\titleformat{\subsection}[block]
  {\large\bfseries}{\thesubsection.}{0.5em}{}

% ----- Tiêu đề "TABLE OF CONTENTS" căn giữa -----
\renewcommand{\contentsname}{TABLE OF CONTENTS}
\titleformat{name=\chapter,numberless}[block]
  {\centering\LARGE\bfseries}{}{0pt}{}
\titlespacing*{name=\chapter,numberless}{0pt}{0pt}{1.5cm}

% ----- Định dạng các dòng trong mục lục -----
\titlecontents{chapter}[0pt]
  {\addvspace{10pt}\bfseries}
  {\thecontentslabel.\ \MakeUppercase}
  {\MakeUppercase}
  {\hfill\contentspage}
\titlecontents{section}[1.5em]
  {\addvspace{2pt}}
  {\thecontentslabel.\ }
  {}
  {\hfill\contentspage}
\titlecontents{subsection}[3em]
  {\addvspace{2pt}}
  {\thecontentslabel.\ }
  {}
  {\hfill\contentspage}


\begin{document}

% ===================== TRANG BÌA =====================
\thispagestyle{empty}

\begin{tikzpicture}[remember picture, overlay]
    \draw[line width=2pt]
        ($(current page.south west) + (0.5in, 0.5in)$)
        rectangle
        ($(current page.north east) - (0.5in, 0.5in)$);
\end{tikzpicture}

\begin{center}
    {\Large \textbf{University of Science and Technology of Hanoi}}\\[0.5cm]

    \includegraphics[width=0.5\textwidth]{images/usth.png}\\[1cm]

    {\LARGE \textbf{Science and Technology Management}}\\[1cm]

    \rule{\textwidth}{0.75pt}\\[0.5cm]
    {\LARGE \textbf{Analysis of a company's technology management activities}}\\[0.5cm]
    \rule{\textwidth}{0.75pt}\\[1.0cm]
\end{center}

% ----- Danh sách thành viên -----
\begin{center}
    \large
    \renewcommand{\arraystretch}{1.1}
    \begin{tabular}{l@{\hspace{1.75cm}}r}
        \multicolumn{2}{l}{\hspace{-1.2cm}\textbf{\Large Members:}} \\[0.3cm]
        \textbf{Nguyen Huy Hai} & 2640088\\
        \textbf{Dang Trung Hieu} & 2640076\\
        \textbf{Pham Mai Hanh} & 2640013\\
        \textbf{Luyen Pham Ngoc Khanh} & 2640014\\
        \textbf{Vuong Hong Minh} & 2640110\\
        \textbf{Trinh Quang Minh} & 2640018\\
        \textbf{Nguyen The Anh} & 2640003\\
        \textbf{Nguyen Thi Minh Anh} & 2640059\\
        \textbf{Nguyen Nam Anh } & 2640058
    \end{tabular}
\end{center}

\vfill
\begin{center}
    \textit{Hanoi, Nov 2026}
\end{center}
\newpage


% ===================== TABLE OF CONTENTS =====================
% \setcounter{page}{2}
\tableofcontents
\listoffigures
% Marks a source that has no entry in references.bib yet.
\providecommand{\pending}[1]{\textbf{[#1]}}

% ===================== MAIN CONTENT =====================
% -------------------------------------------------
% CHAPTER I 
% Responsible: Hai
% executive summary and 1. Company background and the AI-First shift
% -------------------------------------------------
\chapter{Introduction}
% executive summary

\chapter{Company background and the AI-First shift}

\chapter{Technology Identification}

Following Gregory's (1995) technology-management process, technology identification involves scanning the external environment for technologies that may threaten or strengthen a firm's competitive position. For FPT, this process identified artificial intelligence as both a threat to its existing service model and an opportunity to create new forms of value.

\section{Key drivers of AI identification}

\textbf{Market.} Gartner's forecasts, as reported by FPT, put global AI spending at approximately US\$1.5 trillion in 2025 and above US\$2 trillion in 2026 \cite{M3}. This scale of expected demand made AI strategically difficult to ignore.

\textbf{Clients.} Customers were moving from small AI pilots towards company-wide, platform-driven deployment. FPT's large AI transformation engagements illustrate this shift from experimentation towards integration into core operations \cite{A22}.

\textbf{Competition.} AI coding tools can reduce the human effort required for software development, challenging an outsourcing model whose revenue is closely related to engineer headcount and billable man-hours.

\textbf{Policy.} Vietnam's national strategic-technology agenda and Japan's emphasis on sovereign AI capacity increase the value of local infrastructure, trusted implementation and domestic delivery capability \cite{A8,A15}.

\section{Strategic implications for FPT}

The four signals produced three linked implications for FPT's technology strategy.

\begin{table}[H]
\centering
\caption{Strategic implications of AI identification for FPT}
\label{tab:ai-identification}
\begin{tabular}{p{0.19\textwidth} p{0.72\textwidth}}
\toprule
\textbf{Dimension} & \textbf{Implication} \\
\midrule
Threat & AI may reduce demand for headcount-based outsourcing and billable man-hours. \\
Opportunity & FPT can expand into AI consulting, system integration, enterprise applications and the infrastructure required to run them. \\
Response & FPT's AI-First strategy positions AI as a core capability, combining existing models with FPT's platform and integration capabilities. \\
\bottomrule
\end{tabular}
\end{table}

Taken together, these drivers explain why FPT treats AI Factory infrastructure, enterprise platforms and AI-enabled services as strategic priorities. The AI-First response seeks to convert pressure on the traditional delivery model into higher-value consulting, integration and infrastructure services.

% -------------------------------------------------
% Technology Selection and Acquisition 
% Responsible: ?
% -------------------------------------------------

\chapter{Technology Selection}

\section{Technology selection approach}

Technology selection can use scenario analysis, portfolio analysis such as SWOT or BCG, multi-criteria decision-making (MCDM), 
expert judgement and financial analysis. FPT's choice between two discrete strategic alternatives is most clearly represented as 
a multi-attribute decision-making (MADM) problem. The comparison below is our qualitative analysis based on public evidence, not 
FPT's publicly documented selection method.

\section{Comparison of strategic alternatives}

\begin{table}[H]
\centering
\caption{Comparison of strategic alternatives}
\label{tab:technology-selection}
\renewcommand{\arraystretch}{1.15}
\begin{tabular}{p{0.20\textwidth} p{0.36\textwidth} p{0.36\textwidth}}
\toprule
\textbf{Criteria} & \textbf{A. Develop own frontier model} & \textbf{B. Build on existing models} \\
\midrule
Strategic fit & Far from FPT's core IT-services business. & Leverages FPT's existing clients and delivery teams. \\
Feasibility & Requires enormous computing capacity and scarce research talent. & Uses NVIDIA GPUs, partner models and open models. \\
Investment & Mistral raised EUR~1.7bn in 2025 and EUR~3bn in 2026. These are fundraising figures, not directly comparable infrastructure costs \cite{S20,S21}. & Approximately US\$200m per AI Factory site, an illustrative infrastructure figure \cite{A17}. \\
Decision & Not selected. & Selected. \\
\bottomrule
\end{tabular}
\end{table}

\section{Selected technology priorities}

\textbf{AI Factory.} GPU infrastructure for training and running enterprise AI services.

\textbf{FleziPT.} An AI delivery platform launched in August 2025 \cite{S18}.

\textbf{FPT.AI.} AI agents and related enterprise applications.

\textbf{Partner and open models.} Integration of models supplied through partners and open-model ecosystems, allowing FPT to focus on application and delivery capabilities.

Semiconductors remain a longer-term strategic option rather than the central AI selection. FPT made its first commercial shipment of power-management ICs to Japan through Restar on 30 December 2025 \cite{M4,M5,S19}.

\chapter{Technology Acquisition}

\section{Technology acquisition methods}

FPT acquires technology through three routes: collaboration, buying technology or firms, and making technology through internal research and development. Collaboration is the main approach because it combines external assets and networks with FPT's own delivery capability.

\begin{table}[H]
\centering
\caption{FPT's technology acquisition routes}
\label{tab:acquisition-routes}
\renewcommand{\arraystretch}{1.15}
\begin{tabular}{p{0.16\textwidth} p{0.76\textwidth}}
\toprule
\textbf{Route} & \textbf{Evidence} \\
\midrule
Collaborate & NVIDIA strategic partnership (April 2024); Sumitomo and SBI each acquired 20\% of FPT Smart Cloud Japan in 2025, leaving FPT with 60\%; and partnerships with OpenAI and NAVER Cloud in 2026 \cite{S5,S6,A6,A7}. \\
Buy & NVIDIA H100, H200 and HGX B300 GPUs; acquisition of NAC in Japan in 2024 and David Lamm Consulting in Germany in 2025 \cite{A15,A21,S12}. \\
Make & In-house research and development, including US and Taiwan patents, the FleziPT platform, AI Factory design and operation, and proprietary power-management chip designs \cite{S14,S18,S19}. \\
\bottomrule
\end{tabular}
\end{table}

\section{Rationale and trade-offs}

Collaboration dominates because it shares investment and risk, combines complementary capabilities and customer networks, and accelerates entry into new markets. The trade-off is dependence on NVIDIA as the key GPU supplier, alongside dependence on external model partners for part of the technology stack.

\section{Case study: FPT AI Factory}

\begin{table}[H]
\centering
\caption{FPT AI Factory development timeline}
\label{tab:ai-factory-timeline}
\renewcommand{\arraystretch}{1.15}
\begin{tabular}{p{0.14\textwidth} p{0.78\textwidth}}
\toprule
\textbf{Year} & \textbf{Milestones} \\
\midrule
2024 & NVIDIA partnership in April; AI Factory launches in Japan and Vietnam in November \cite{S5,S7}. \\
2025 & Sumitomo and SBI acquired a combined 40\% of the Japanese unit. In June, the sites ranked Japan 36th (49.85 PFLOPS) and Vietnam 38th (46.65 PFLOPS) in TOP500 \cite{S6,S10}. \\
2026 & NVIDIA HGX B300 announced in March. Average utilisation was approximately 60\% in February, with Vietnam at breakeven and Japan loss-making. Both sites were later reported above 90\% utilisation and profitable in Q2 \cite{A21,S4}. \\
\bottomrule
\end{tabular}
\end{table}

\begin{figure}[H]
\centering
\begin{tikzpicture}[x=0.095\textwidth,y=0.65cm,font=\small]
\draw[gray] (0,-0.45) -- (0,1.45);
\foreach \y/\v/\label in {1/60/{February 2026},0/90/{Q2 2026}}{
  \node[anchor=east] at (0,\y) {\label};
  \fill[blue!55!black] (0,\y-0.22) rectangle (\v,\y+0.22);
  \node[anchor=west] at (\v,\y) {\ifnum\v=90 $>90\%$\else $\approx60\%$\fi};
}
\end{tikzpicture}
\caption[AI Factory utilisation]{Reported AI Factory utilisation in February and Q2 2026. Sources: FPT and DNSE \cite{S4}.}
\label{fig:ai-factory-utilisation}
\end{figure}

AI Factory revenue was estimated at approximately US\$12~million in H1 2026 \cite{S4}. The utilisation trend and reported site-level results indicate operational improvement and commercial progress. Since separate profit-and-loss statements are unavailable, site-level profitability rests on company disclosure. It does not establish profitability or ROI for FPT's overall AI business.



% -------------------------------------------------
% CHAPTER III - Technology Exploitation & Commercialisation
% Responsible: Minh
% at 5. Technology exploitation and commercialisation
% -------------------------------------------------
\chapter{Technology Exploitation \& Commercialisation}

A technology creates value only when customers adopt and pay for it. FPT earns from AI in three ways, summarised in Table~\ref{tab:channels}.

\section{Three commercialisation channels}

\begin{table}[htbp]
  \centering
  \small
  \caption{How FPT earns from AI}
  \label{tab:channels}
  \begin{tabular}{>{\raggedright\arraybackslash}p{2.8cm}>{\raggedright\arraybackslash}p{4.6cm}>{\raggedright\arraybackslash}p{6.4cm}}
    \toprule
    \textbf{Channel} & \textbf{What clients pay for} & \textbf{Real-world example} \\
    \midrule
    AI transformation projects & Design and build of AI systems across a client's business & Contract worth tens of millions of US dollars with a global materials group (April 2026) \cite{S25} \\ \addlinespace
    AI agents & Ready-made AI agents for uses such as customer service, sales and document processing & 10,220 FPT.AI agents deployed (2025) \cite{A26} \\ \addlinespace
    GPU cloud & Rented computing capacity for training and running AI models & Analysts report that AI Factory clients are mainly banks and technology companies \cite{M6} \\
    \bottomrule
  \end{tabular}
\end{table}

\textbf{AI transformation projects:} In April 2026, FPT signed a contract worth tens of millions of US dollars with a global materials group to move the group toward an agentic AI operating model on FleziPT \cite{A22}. In May 2026, it announced 6 AI transformation partnerships in Thailand and Singapore, with partners including UOB, SCG, SATS and ComfortDelGro; these include memoranda of understanding and are not described as paid contracts \cite{S25}. In 2025, it signed a record USD~256~million AI-powered digital solutions contract with an energy company \cite{A27}.

\textbf{AI agents:} In 2025, FPT.AI reported 10,220 AI agents deployed in business operations and more than 200 customers in 15 countries. Separately, FPT's language models processed 169 million requests \cite{A26}. Typical uses include customer service, sales support and document processing.

\textbf{GPU cloud:} Companies rent AI Factory capacity to train or run their own models. Analysts report that these clients are mainly banks and technology companies \cite{M6}.

\section{Financial results}

AI and data analytics revenue was VND~2,612 billion in 2025 (+40.7\%), equivalent to about USD~100~million and 15.5\% of digital-transformation revenue \cite{M1,M6}. It reached VND~1,842 billion in H1 2026 (+54\%) and VND~2,915 billion in 9M2026 (+55\%), about 7.3\% of group revenue (our calculation: 2,915 / 40,023) \cite{M14,S2}.

At foreign level, FPT won 16 projects worth more than USD~10~million each in 9M2026, down from 19 a year earlier, while the number of contracts of at least USD~1~million rose 29\% to 163 \cite{M15}. In 2025, newly signed AI and cloud revenue was VND~1,540 billion, and recognised revenue was VND~913 billion (+48.4\%) \cite{S16}.

These figures should be interpreted with caution. Adoption figures such as agents, requests show usage, not profit. Judging commercial success requires revenue, margin and client retention data, and FPT discloses only part of this.




% -------------------------------------------------
% CHAPTER IV - Technology Protection, Governance & Risk
% Responsible: Minh
% at 6. Technology protection, governance and risk
% -------------------------------------------------
\chapter{Technology Protection, Governance \& Risk}
Protection means keeping control of the value a technology creates. For FPT, this has three parts: intellectual property, certified governance and regulatory compliance.

\section{Intellectual property}

FPT protects its intellectual property through the patents described in Chapter~II, through proprietary platform know-how, and through contracts with clients and partners. Because the base models belong to OpenAI, NAVER or open-source projects, FPT's patents protect only its own application layers. In practice, process know-how and client relationships are more important protections.

\section{Certified governance}

FPT's governance is evidenced by external certification:

\begin{itemize}
  \item \textbf{ISO/IEC 42001 (AI management).} FPT Software was the first Vietnamese company to receive accredited certification for ISO/IEC 42001, the international standard for AI management systems \cite{S23,M11}.
  \item \textbf{SOC~2 and SOC~3 (cloud security).} The Vietnam AI Factory holds SOC~2 Type~II and SOC~3 reports, while the Japan site holds a SOC~2 Type~I report \cite{S28}.
  \item \textbf{HITRUST AI Security (AI security).} In February 2026, FPT's Hanoi data centre received HITRUST r2 certification, including the HITRUST AI Security Certification, which covers threats such as data poisoning, model inversion and prompt injection \cite{M18}.
\end{itemize}

\section{Regulatory compliance: Vietnam's AI Law}

Regulation is now a formal constraint. Vietnam's Law on Artificial Intelligence (No.~134/2025/QH15) was passed on 10~December 2025 with 429 of 434 votes and took effect on 1~March 2026 \cite{M19}. The law:

\begin{itemize}
  \item classifies AI systems into three risk tiers;
  \item requires conformity assessment for high-risk systems;
  \item requires important AI applications in essential sectors to run on national AI infrastructure;
  \item gives existing systems in the medical, education and financial sectors 18 months to comply \cite{M20}.
\end{itemize}

The law adds compliance costs to FPT's banking and healthcare work. At the same time, its emphasis on local infrastructure could favour domestic providers; this second point is our interpretation, not a benefit stated by FPT.

\section{Key risks and management responses}

Table~\ref{tab:risks} sets out the key risks identified in our analysis, together with possible management responses. These are analytical risks, not failures that FPT has reported.

\begin{table}[htbp]
  \centering
  \small
  \caption{Key risks and possible management responses}
  \label{tab:risks}
  \begin{tabular}{>{\raggedright\arraybackslash}p{2.8cm}>{\raggedright\arraybackslash}p{6.4cm}>{\raggedright\arraybackslash}p{4.6cm}}
    \toprule
    \textbf{Risk} & \textbf{Description} & \textbf{Possible response} \\
    \midrule
    Idle GPUs & GPU costs and idle capacity if demand weakens & Utilisation dashboards and pre-sold capacity \\ \addlinespace
    Dependence on key suppliers & Dependence on NVIDIA for chips, and on OpenAI and NAVER for models & Multi-vendor and multi-model architecture \\ \addlinespace
    AI errors and privacy & Model errors and data privacy incidents at clients & Human review, audit logs and systematic model evaluation \\ \addlinespace
    Fewer billable hours & Revenue cannibalisation, as faster AI delivery reduces billable effort & Outcome-based pricing \\
    \bottomrule
  \end{tabular}
\end{table}


% -------------------------------------------------
% CHAPTER V - Performance Evaluation 
% Responsible: Khanh
% at 7. Performance evaluation
% -------------------------------------------------
\chapter{Performance Evaluation}

\section{Investment and delivery capability}

We assess investment as moderate relative to developing a frontier foundation model. FPT builds on existing models and shares the Japanese AI Factory investment with minority shareholders Sumitomo and SBI \cite{A5,A8}. Co-investment limits the capital FPT must supply alone. In Q2 2026, FPT reported utilisation above 90\% and a profit at each of its Vietnamese and Japanese clusters \cite{A29}. This shows demand for installed capacity. One profitable quarter at two sites is insufficient to judge the efficiency of the investment over its life.

The THIO framework helps explain FPT's delivery capability. Its GPU infrastructure (technoware) complements staff expertise developed through research collaboration with Mila (humanware) \cite{A16,A13}. FleziPT captures engineering knowledge in a reusable platform (infoware), while the certified AI management system gives teams common procedures (orgaware) \cite{A22,S23}. Together these support our assessment of strong acquisitive and operative capabilities: FPT can obtain technology through partners and run it at scale. Its innovative capability is applied, as illustrated by its granted US patent for acoustic anomaly detection \cite{S13}. Reliable component ratings are unavailable, so we do not calculate a Technology Contribution Coefficient.

\section{Adoption and profitability}

In-house delivery is the main commercial route. The five-year US\$256m energy-sector contract signed in September 2025 commits FPT to deliver AI-augmented services. Its face value covers the entire engagement \cite{S24}. The six AI transformation memoranda of understanding signed in 2026, including those with UOB, SCG, SATS and ComfortDelGro, indicate prospective demand. Signing them establishes neither deployment nor recognised revenue \cite{S25}. Evidence of actual use comes from FPT.AI, which reported more than 200 enterprise clients and 10,220 active agents in 2025 \cite{A26}, and from FPT's reported deployment of Flezi Metis for an automotive technology group \pending{FPT 9M2026 results}.

Joint commercialisation extends this route through the Japanese venture and the collaboration with NAVER Cloud on AI Factory and GPU-as-a-Service offerings \cite{A5,A7}. The sources reviewed provide no public evidence of a separate route based on selling or licensing FPT's AI technology.

\begin{figure}[H]
\centering
\begin{tikzpicture}[x=0.11cm,y=0.6cm,font=\small]
\draw[gray] (0,-0.45) -- (0,1.45);
\foreach \y/\v/\name in {1/12.4/Whole group,0/55/AI and data analytics}{
  \node[anchor=east] at (0,\y) {\name};
  \fill[blue!55!black] (0,\y-0.23) rectangle (\v,\y+0.23);
  \node[anchor=west] at (\v,\y) {+\v\%};
}
\end{tikzpicture}
\caption[Revenue growth in the first nine months of 2026]{Year-on-year revenue growth, January--September 2026. Sources: \href{https://fpt.com/en/news/fpt-s-9-months-2026-revenue-and-pbt-increase-12-4-and-17-0}{FPT nine-month results} and CafeF \cite{S2}.}
\label{fig:ai-growth}
\end{figure}



\section{Revenue contribution and strategic progress}

AI and data analytics generated VND~2,915bn in January--September 2026, already above the VND~2,612bn recorded in the whole of 2025 \cite{M1,S2}. On a year-on-year basis, revenue in this category grew by 55\%, against 12.4\% for the group after restating the prior period for FPT Telecom (Figure~\ref{fig:ai-growth}) \pending{FPT 9M2026 results}. The 42.6 percentage point gap shows a shift in the revenue mix from a small base. Dividing VND~2,915bn by group revenue of VND~40,023bn gives 7.3\% for the nine-month period. This measures the disclosed services category and cannot isolate AI embedded in other contracts.

Shared investment, sustained customer use and recognised revenue support a qualitative judgement that FPT's strategy is coherent and operationally on track. The available results do not disclose an operating margin or return on investment (ROI) for the AI business as a whole. Its overall financial payoff therefore remains unproven.


% -------------------------------------------------
% CHAPTER VI - CONCLUSION 
% Responsible: Khanh (delete)
% at 8. Conclusion
% ------------------------------------------------
\chapter{Conclusion}

Gregory's five processes explain the logic of FPT's strategy. In identification and selection, AI posed a threat to billable software work while creating demand for enterprise automation. FPT chose to build solutions on existing models and acquired much of the necessary capability through collaboration \cite{A8,A13,A16}. It exploits that capability mainly through its own delivery teams, supplemented by joint commercialisation. Revenue in the reported AI and analytics category is growing rapidly but still makes up a small share of group sales \pending{FPT 9M2026 results}. Protection centres on what FPT develops: proprietary platforms, implementation know-how and granted patents for applications such as acoustic anomaly detection \cite{S13}. Its certified management system supports oversight, while NVIDIA GPUs and third-party models remain external technologies on which delivery depends \cite{S23,A21,A7}.

Existing customer relationships help FPT identify work worth automating. Its delivery scale lets it serve clients across markets, and its technical and industry knowledge helps teams adapt models to specific workflows. These advantages support a coherent strategy and credible operational progress. The financial payoff remains open. Future assessments should track AI operating margins and ROI, exposure to GPU and model suppliers, and AI's share of group revenue under a consistent definition.


% ===================== REFERENCES =====================
\chapter*{References}
\addcontentsline{toc}{chapter}{References}
\printbibliography[heading=none]


\end{document}